% Track d5 -- Duplication: the mathematics of Cl(3,3) -> Cl(6,6).
% Verification: investigations/rest_energy/verify/d5_duplication_math.py (claims d5-1a ... d5-10).
% New references are in sections/references.tex.
\section{\texorpdfstring{Duplication: from $\Cl(3,3)$ to $\Cl(6,6)$}{Duplication: from Cl(3,3) to Cl(6,6)}}
\label{d5:sec:main}

The author's second remark asks which physical and mathematical process makes a unitary energy ``result in a duplication'', from
$\Cl(3,3)$ to $\Cl(6,6)$. A duplication acts on four levels at once: on the vector space, $V\to V_2:=V\oplus V$; on the algebra,
$\Cl(V)\to\Cl(V_2)=\Cl(6,6)$; on the spinor module, $S\to S_2=S\otimes S$; and on the energy, $E_6\to E_6\oplus E_6$. This section takes
the four steps in turn, asks which are forced and which are chosen, and records what the duplicate can do that one copy cannot. The
answers, with their classes, are these. The step on the algebra is forced: it is the graded tensor product, its Klein factor is the
Jordan--Wigner sign, and $S_2$ is the Fock space of six fermionic modes (inherited; \cref{d5:sec:algebra}). The step on the energy is a
choice within a three-sphere of energy-preserving, Lorentz-covariant duplications, which is the algebra of a lossless beam splitter. Only a
circle in it, the copy rotation, respects $\eta$, and the two notions of rest energy that coincide on one copy separate on two (fixed;
\cref{d5:sec:energy}). The duplicate carries structures one copy cannot carry: an $\eta$-isometric complex structure, a compact,
Lorentz-invariant copy charge, the copy algebra $\so{2,\C}$, a closure that runs at constant rate with non-null copies ($E_0(\pi^{(k)})\neq0$), and a non-rigid,
Lorentz-invariant vacuum (fixed, partly restated from \cite{ES17}; \cref{d5:sec:closure,d5:sec:actions,d5:sec:spin,d5:sec:vacuum}). Finally, the
copy rotation is the only admissible $b$-field between two binding frames that commutes with the binding root and the rotations. When
the two rates differ it cannot be removed, and it confines a frame that has no rate of its own; it creates no rate where both frames
are free (fixed, with the mixing formulas inherited; \cref{d5:sec:bfield}). Every finite-dimensional statement was checked by
\texttt{verify/d5\_duplication\_math.py}; identifiers d5-$n$ refer to its \texttt{PASS} lines (\cref{d5:tab:claims}).

\paragraph{Notation.} $V_2=V\oplus V$ has the basis $e_1,\dots,e_6$ of copy~$1$ followed by that of copy~$2$, the metric
$\eta_2=\id_2\otimes\eta$, and the binding root $G_2=\id_2\otimes G$ of both copies (as $G_3$ in Theorem~17.19 of \cite{ES17}).
$\varepsilon=\bigl(\begin{smallmatrix}0&-1\\1&0\end{smallmatrix}\bigr)$, $\sigma_x$, $\sigma_z$ act on the copy index, and
$C_v:=\varepsilon\otimes\id_6$ is the copy rotation on vectors. A pair of field strengths is $(z_1,z_2)\in\C^3\oplus\C^3$, with
$q_2:=q(z_1)+q(z_2)$ and duplicated energy $E_6(z_1)+E_6(z_2)$. On $S_2=\R^8\otimes\R^8$ the tower of equation~(2) of \cite{ES17} gives
$e^{(1)}_a=e_a\otimes\id$, $e^{(2)}_a=\omega\otimes e_a$, $\omega=e_1\cdots e_6$, and $E_{6(k-1)+a}=e^{(k)}_a$. Statements about $V_2$ are
pointwise linear algebra and apply to whichever field is placed there. Where the field matters we name it: the frame field $X$, or the
momentum field $\pi=\eta^{-1}p$ of the closure, to which the Casimirs belong (Section~17 of \cite{ES17}, \cref{d3:sec:main}). The
rates $\kappa_k$ of the binding, the spectrum of the self-adjoint part of the shift, are kept apart from the lines $\omega$ of the motion, the
spectrum of the shift itself: a $b$-field can move the lines but never the rates (\cref{d7:thm:schr}(iii)). Reading either as a rest energy through
$mc^2=\hbar\kappa$ is an added identification.

\subsection{The duplicated algebra: graded tensor product and Klein factor}
\label{d5:sec:algebra}

\begin{proposition}[The duplicated algebra]
\label{d5:prop:graded}
\begin{enumerate}[label=(\roman*),leftmargin=2.2em]
\item $\omega^2=\id$, $\omega^{\!\top}=\omega$, $\Tr\omega=0$, and $\omega$ anticommutes with every $e_a$. The twelve matrices $e^{(k)}_a$ satisfy
the Clifford relations of $\R^{6,6}$. Their $4096$ blades are signed permutation matrices, and all except $\id$ are traceless, so
$\Cl(6,6)\cong M_{64}(\R)$.
\item For blades $e_I,e_J$ of $\Cl(3,3)$, $e^{(1)}_Ie^{(2)}_J=(e_I\,\omega^{|J|})\otimes e_J$ and $e^{(2)}_Je^{(1)}_I=(-1)^{|I||J|}e^{(1)}_Ie^{(2)}_J$.
Hence $a\,\hat\otimes\,b\mapsto a\,\omega^{|b|}\otimes b$ is an algebra isomorphism $\Cl(3,3)\hat\otimes\Cl(3,3)\to\Cl(6,6)$, where $\hat\otimes$
is the graded tensor product, $(a\hat\otimes b)(c\hat\otimes d)=(-1)^{|b||c|}ac\,\hat\otimes\,bd$.
\item The Klein factor is forced. In $M_{64}(\R)$ the elements commuting with every $e_a\otimes\id$ form $\id\otimes M_8(\R)$, and those
anticommuting with every $e_a\otimes\id$ form $\omega\otimes M_8(\R)$. Any realisation of $\Cl(V_2)$ on $\R^{64}$ that extends
$e_a\mapsto e_a\otimes\id$ is therefore conjugate to the tower by an element of $\id\otimes GL(8,\R)$.
\end{enumerate}
\end{proposition}

\begin{proof}
(i) $\omega^2=(-1)^{15}(-1)^3\id$ by the transpose rule (4) of \cite{ES17}, and $\omega$ has even grade in an even dimension. The relations
follow from $(\omega\otimes e_a)(\omega\otimes e_b)=\id\otimes e_ae_b$ and $(e_a\otimes\id)(\omega\otimes e_b)=-(\omega\otimes e_b)(e_a\otimes\id)$. A product
of signed permutation matrices is one, and if $I\neq J$ then $e_I^{\!\top}e_J=\pm e_{I\triangle J}$ is traceless, so the blades are orthogonal
for the trace form; $4096=64^2$. (ii) $e^{(2)}_J=\omega^{|J|}\otimes e_J$, and moving the $|J|$ copy-2 generators past the $|I|$ copy-1
generators costs $(-1)^{|I||J|}$. (iii) By (3) of \cite{ES17} a blade $e^{(1)}_Ie^{(2)}_J$ commutes with $e^{(1)}_a$ exactly when
$|I|+|J|-[a\in I]$ is even. For all six $a$ this forces either $I=\emptyset$ and $|J|$ even, or $I=\{1,\dots,6\}$ and $|J|$ odd: $64$ blades,
each $\pm\id\otimes e_J$. The anticommuting case is the same with the parities reversed. If $e'_a$ anticommutes with all $e_b\otimes\id$, then
$e'_a=\omega\otimes f_a$, and $(\omega\otimes f_a)(\omega\otimes f_b)=\id\otimes f_af_b$, so the $f_a$ satisfy the relations of $\Cl(3,3)\cong M_8(\R)$. By the
Skolem--Noether theorem $f_a=ue_au^{-1}$, and conjugation by $\id\otimes u$ fixes $e_a\otimes\id$ and $\omega\otimes\id$.
\end{proof}

Everything here is inherited: it is the exponential law $\Cl(V\oplus W)\cong\Cl(V)\hat\otimes\Cl(W)$ \cite{Chevalley54,LawsonMichelsohn89}, and
(i) restates equation~(2) of \cite{ES17}. The arena contributes only the fixed realisation (d5-1a to d5-1d). The duplication of the
algebra is not a choice; given copy~$1$, copy~$2$ is determined up to conjugation.

\begin{proposition}[Fock space and Jordan--Wigner strings]
\label{d5:prop:fock}
Put $n^\pm_i=\tfrac12(e_i\pm e_{i+3})$, $i=1,2,3$. Then $\{n^\pm_i,n^\pm_j\}=0$ and $\{n^+_i,n^-_j\}=\delta_{ij}$, so $\Cl(3,3)$ is the real
algebra of canonical anticommutation relations of three fermionic modes and $S=\R^8$ is their Fock space. The occupation numbers are
$N_i=n^+_in^-_i=\tfrac12(\id-\omega_i)$, with $\omega_i=e_ie_{i+3}$ the pole operators. The eight eigenlines of Proposition~2.5(i) of
\cite{ES17} are therefore the eight Fock states, and $(-1)^{N_1+N_2+N_3}=-\omega$. At $n=2$ the tower supplies six modes. The modes of copy~$2$ are
$\omega\otimes n^\pm_i$, each carrying the string $-(-1)^{N^{(1)}}$ of copy~$1$. The total number has spectrum $0,\dots,6$ with multiplicities
$1,6,15,20,15,6,1$, and $(-1)^N$ is the volume element of $\Cl(6,6)$.
\end{proposition}

Inherited. In split signature the vectors $e_i\pm e_{i+3}$ span two maximal isotropic subspaces, and the spinor module is the exterior
algebra of one of them \cite{Chevalley54}. The string is the Jordan--Wigner transformation \cite{JordanWigner28}, and the Klein factor is the
device that makes the fields of two subsystems anticommute \cite{Klein38}. On the level of states, the mathematical process behind
``$\Cl(3,3)\to\Cl(6,6)$'' is the second-quantisation rule $\mathrm{Fock}(\mathcal H\oplus\mathcal H)=\mathrm{Fock}(\mathcal H)\otimes\mathrm{Fock}(\mathcal H)$
(d5-1e).

\subsection{Centre-of-mass and relative halves}
\label{d5:sec:cm}

\begin{proposition}[Two halves, and the exchange]
\label{d5:prop:cmrel}
Let $f_a=(e^{(1)}_a+e^{(2)}_a)/\sqrt2$ and $g_a=(e^{(1)}_a-e^{(2)}_a)/\sqrt2$.
\begin{enumerate}[label=(\roman*),leftmargin=2.2em]
\item $\{f_a,f_b\}=\{g_a,g_b\}=2\eta_{ab}$ and $\{f_a,g_b\}=0$ for all $a,b$, so $\Cl(6,6)=\Cl(f)\hat\otimes\Cl(g)$: two mutually anticommuting
copies of $\Cl(3,3)$.
\item The copy rotation $C_b=\tfrac12\sum_a\eta_{aa}e^{(1)}_ae^{(2)}_a$ acts on vectors by $e^{(1)}_a\mapsto-e^{(2)}_a$, $e^{(2)}_a\mapsto e^{(1)}_a$, that is by $-C_v$, so it lifts the copy rotation with the opposite orientation; it is a sum of
six commuting terms of square $-\tfrac14$, so $\exp(2\pi C_b)=\id$. Conjugation by $\exp(\tfrac\pi4C_b)$ carries $e^{(1)}_a$ to $g_a$ and $e^{(2)}_a$ to
$f_a$, and $C_b=-\tfrac12\sum_a\eta_{aa}f_ag_a$. The split into halves is the split into constituents turned by an eighth of a turn of the
copy rotation, which rotates one half into the other.
\item Let $\omega=(-1)^{\mathsf p}$ on the basis of equation~(1) of \cite{ES17}, in which $\omega$ is diagonal, let $\tau(x\otimes y)=y\otimes x$, and let
$D=(-1)^{\mathsf p\otimes\mathsf p}$. The Koszul swap $\tau_K=D\tau$ satisfies $\tau_Ke^{(1)}_a\tau_K=e^{(2)}_a$, $\tau_K^2=\id$, $\Tr\tau_K=0$, and
$\tau_K=g_1g_2\cdots g_6$. So the exchange commutes with every $f_a$ and anticommutes with every $g_a$. The naive swap $\tau$ does not implement
the exchange.
\item The summed clock is $e^{(1)}_1+e^{(2)}_1=\sqrt2f_1$. Every diagonal bivector satisfies $\sum_ke^{(k)}_ae^{(k)}_b=f_af_b+g_ag_b$; this covers the
diagonal interior and binding Lorentz generators and the summed charge $J^{(1)}+J^{(2)}=f_2f_3+g_2g_3$, which act on both halves alike
with no cross term. $G_2$ preserves both halves. The copy rotations $C_b$ and $M_{12}$ (the copy rotation of the interior four-planes, Section~10.1 of \cite{ES17}) are pure
cross terms and anticommute with $\tau_K$.
\end{enumerate}
\end{proposition}

\begin{proof}
(i) and (iv) are expansions, using $e^{(1)}_ae^{(2)}_a=-f_ag_a$. (ii) By (3) of \cite{ES17}, $[C_b,e^{(1)}_b]=-e^{(2)}_b$. The six terms
$\tfrac12\eta_{aa}e^{(1)}_ae^{(2)}_a$ share no index, so they commute, and each squares to $-\tfrac14$; hence
$\exp(\theta C_b)=\prod_a(\cos\tfrac\theta2+\sin\tfrac\theta2\,\eta_{aa}e^{(1)}_ae^{(2)}_a)$, which is $(-1)^6\id$ at $\theta=2\pi$. On vectors $\exp(\theta\ad C_b)$ turns
$e^{(1)}_a$ into $\cos\theta\,e^{(1)}_a-\sin\theta\,e^{(2)}_a$, which is $g_a$ at $\theta=\tfrac\pi4$. (iii) The swap of the copies fixes every $f_a$ and negates every $g_a$: it is the reflection of $V_2$
in the span of the $f_a$ along the span of the $g_a$, and the product $g_1\cdots g_6$ of an even number of unit vectors implements it by conjugation. The Koszul swap implements
it too, by the sign rule. The commutant of $\Cl(6,6)=M_{64}(\R)$ is $\R$, so the two agree up to a scalar, which is $+1$. $\Tr\tau_K$ counts
$\pm1$ over the diagonal states $x_i\otimes x_i$ with sign $(-1)^{\mathsf p_i}$, and $\omega$ is traceless.
\end{proof}

Inherited: centre-of-mass and relative coordinates, and the graded swap, which is the exchange operator of two fermionic subsystems. The
identity $\tau_K=g_1\cdots g_6$, with its sign, and the list in (iv) are facts of the realisation (d5-1f, d5-2a to d5-2c). Because $\tau_K$
anticommutes with the copy rotation, exchanging the constituents reverses the copy charge, as exchanging two real vector fields $W^1,W^2$
exchanges $W^\pm$ (inherited). The halves are not invariant under the copy rotation. Only the exchange selects them, as its two eigenspaces
on vectors.

\subsection{Duplicating the undivided energy}
\label{d5:sec:energy}

\begin{proposition}[Symmetry of the duplicated energy]
\label{d5:prop:energysym}
\begin{enumerate}[label=(\roman*),leftmargin=2.2em]
\item The linear maps of $V_2$ that commute with the diagonal rotations and with $G_2$ form $\mathfrak{gl}(2)\otimes\operatorname{span}\{\id,G\}\cong\mathfrak{gl}(2,\C)$.
The infinitesimal symmetries of $E_6\oplus E_6$ among them form $\uu2=\operatorname{span}\{C_v,\ \id\otimes G,\ \sigma_z\otimes G,\ \sigma_x\otimes G\}$. Those
in $\so{6,6}$ form $\so{2,\C}=\operatorname{span}\{C_v,\ \varepsilon\otimes G\}$, and those in both form $\so2=\R\,C_v$. The overall duality $\id\otimes G=G_2$, the
relative duality $\sigma_z\otimes G$ (the binding rotation in opposite senses on the two copies) and $\sigma_x\otimes G$ are $\eta_2$-self-adjoint,
like $G$.
\item Inside $\uu6=\so{12}\cap\Cent(G_2)$, of dimension $36$, the diagonal $\uu3$ and $\uu2\otimes\id$ are each other's commutants: $U(3)\times U(2)$ is
a dual pair in $U(6)$.
\end{enumerate}
\end{proposition}

Part (i) restates rows three to five of Proposition~13.14 of \cite{ES17} (fixed there), with the generators named. Row three is stated there for the diagonal $\lG$, whose commutant in $\so{6,6}$ is the same $\so{2,\C}$ (\cref{d5:prop:so6C}(iii)). The four generators are
those of a lossless beam splitter on two modes \cite{CamposSalehTeich89}: a real splitter $C_v$, an overall phase, a relative phase and a
symmetric splitter. Their group $U(2)$ is the energy-preserving duality group of two free Maxwell fields \cite{GaillardZumino81}, as
Section~13.5 of \cite{ES17} notes. Part (ii) is inherited from Howe duality for $\C^3\otimes\C^2$ \cite{Howe89} (d5-3a, d5-3b).

\begin{theorem}[Duplication maps]
\label{d5:thm:dup}
\begin{enumerate}[label=(\roman*),leftmargin=2.2em]
\item The $G$-linear, rotation-equivariant maps $V\to V_2$ are $\iota_{\alpha,\beta}(z)=(\alpha z,\beta z)$ with $\alpha,\beta\in\C$. Such a map preserves
the energy, $E_6(\alpha z)+E_6(\beta z)=E_6(z)$, if and only if $|\alpha|^2+|\beta|^2=1$, which gives a three-sphere $U(2)/U(1)$: the first column of
a beam splitter. It is an isometry $(V,\eta)\to(V_2,\eta_2)$ if and only if $\alpha^2+\beta^2=1$ over $\C$, a non-compact complex conic, the orbit
of $\iota_{1,0}$ under the group of the copy algebra $\so{2,\C}$. An energy-preserving $\iota_{\alpha,\beta}$ is an isometry if and only if
$\alpha,\beta$ are real with $\alpha^2+\beta^2=1$, which gives a circle, the orbit of the copy rotation.
\item On the image, $q_2=(\alpha^2+\beta^2)\,q$, $E_0^{(1)}+E_0^{(2)}=(|\alpha|^2+|\beta|^2)E_0$, and the summed Maxwell stress tensor is
$T(\alpha z)+T(\beta z)=(|\alpha|^2+|\beta|^2)\,T(z)$. Every energy-preserving duplication therefore preserves the summed stress tensor, the sum
of the copies' rest energies, and the pair invariant $E_{\rm pair}^2:=\tfrac14T_{\rm tot}{}^\mu{}_\nu T_{\rm tot}{}^\nu{}_\mu$, which equals $E_0^2$ by
Rainich's identity $(T\eta)^2=E_0^2\,\id$ \cite{Rainich1925}. By contrast $|q_2|=|\alpha^2+\beta^2|\,|q|=\bigl(1-4(\Im\alpha\bar\beta)^2\bigr)^{1/2}|q|$ sweeps $[0,|q|]$. The complex invariant $q_2$ is preserved exactly on
the circle. Its modulus $|q_2|$ is preserved exactly where $\Im\alpha\bar\beta=0$, that is on the torus $(\alpha,\beta)=e^{i\psi}(\cos t,\sin t)$, the orbit of the copy
rotation and the overall duality rotation $e^{\psi G_2}$; for example $\iota_{i,0}$ gives $q_2=-q$. At the equal split with relative duality angle $\theta$ ($\beta=e^{i\theta}\alpha$), $|q_2|=|\cos\theta|\,|q|$, which vanishes at
$\theta=\pi/2$.
\item Without equivariance, consider $X\mapsto(\cos\phi\,X,\ \sin\phi\,RX)$ with $R\in U(3)$ and $0<\phi<\pi/2$. It preserves the summed stress tensor
for all $X$ iff $R=e^{\vartheta G}$ is a duality rotation, preserves $q_2$ iff $R\in O(3)$, and preserves both iff $R=\pm\id$. More generally,
a $G$-linear embedding $W=U+iV:\C^3\to\C^6$ that preserves both $E_6\oplus E_6$ and $q_2$ is real, since
$W^*W-W^{\!\top}W=2V^{\!\top}V-2iV^{\!\top}U$.
\end{enumerate}
\end{theorem}

\begin{proof}
(i) The rotations act on $V=\R^2\otimes\R^3$ through the absolutely irreducible second factor, so by Schur's lemma \cite{Schur} the equivariant maps
are $(A,B)\otimes\id_3$ with $A,B\in M_2(\R)$, and $G$-linearity puts $A,B$ in $\operatorname{span}\{\id,G\}\cong\C$. With $G^{\!\top}G=\id$,
$G^{\!\top}\eta G=-\eta$ and $G^{\!\top}\eta=\eta G$, the map $L=\iota_{\alpha,\beta}$ has $L^{\!\top}L=(|\alpha|^2+|\beta|^2)\id$ and
$L^{\!\top}\eta_2L=\Re(\alpha^2+\beta^2)\,\eta+\Im(\alpha^2+\beta^2)\,\eta G$. Since $|\alpha|^2+|\beta|^2-\Re(\alpha^2+\beta^2)=2[(\Im\alpha)^2+(\Im\beta)^2]$, the two
conditions together force real coefficients. (ii) $q$ is holomorphic and quadratic. The components of $T$ are $\tfrac12|z|^2$,
the textbook Poynting vector $-\mathbf S=-\tfrac i2\bar z\times z$, and $\tfrac12\delta_{ij}|z|^2-\Re(\bar z_iz_j)$, all sesquilinear, so $T(\alpha z)=|\alpha|^2T(z)$. Rainich's identity is the
tensor form of the identity $E_6^2=E_0^2+|\mathbf S|^2$ of \cref{d1:sec:identity}. (iii) $T$ determines the Hermitian matrix $zz^*$, the real part from $T^{ij}$ and
the imaginary part from $T^{0i}$, and conversely. So $T(Rz)=T(z)$ for all $z$ iff $R$ commutes with every rank-one Hermitian matrix, iff $R$ is a
phase. $q(Rz)=q(z)$ for all $z$ iff $R^{\!\top}R=\id$, which for unitary $R$ means $\bar R=R$. If $W^*W=W^{\!\top}W=\id$ the identity gives
$V^{\!\top}V=0$.
\end{proof}

All three parts are inherited: they are the linear algebra of a beam splitter acting on $\C^3\otimes\C^2$ (d5-3c to d5-3e). What is not
inherited is the reading.

\begin{remark}[Two rest energies of a duplicate]
\label{d5:rem:twoE0}
On one copy $E_0=\tfrac12|q|$ is at once the holomorphic invariant and the stress-energy invariant; this is Rainich's identity, the field-point
form of $E^2=(pc)^2+(mc^2)^2$ (\cref{d1:sec:identity}). On two copies the two separate. The stress-energy content ($T_{\rm tot}$,
$\sum_kE_0^{(k)}$, $E_{\rm pair}$) is preserved by the whole three-sphere, the holomorphic invariant $q_2$ only by the circle (fixed). Reading this as the
point where duplication changes the meaning of ``rest energy'' is interpretive. Which one deserves the name depends on the group. The invariants of a pair
under the diagonal $\so{3,\C}$ are the entries of the complex Gram matrix $Q_{kl}=z_k\cdot z_l$; the copy algebra $\so{2,\C}$ preserves
$\Tr Q=q_2$ and $\det Q$ but not $Q_{11}$ (d5-3f); and the closure sees $|q_2|$ of the momentum field (\cref{d5:prop:balance}). Two
counter-propagating null fields, $z_1=(1,i,0)$ and $z_2=(1,-i,0)$, have $E_0^{(k)}=0$ and $q_2=0$, yet $E_{\rm pair}^2=2$ (d5-3i). Their rest energy as a pair is the invariant mass $M_{\rm pair}=2$ of \cref{d6:thm:pair}, not
$E_{\rm pair}=\sqrt2$: $E_{\rm pair}$ is positive exactly when $M_{\rm pair}$ is, and equals it only when $\det Q=0$, as on every duplicate of one field
(\cref{d6:thm:pair}). A nonzero $M_{\rm pair}$ is the field-point analogue of the invariant mass of two photons, which permits $\gamma\gamma\to e^+e^-$ \cite{BreitWheeler34}. A single null field
duplicated by any $\iota_{\alpha,\beta}$ stays null in each copy, since $q(\alpha z)=\alpha^2q(z)$; that is the analogue of the absence of
$\gamma\to e^+e^-$ in vacuum (the photon kinematics inherited; both analogies interpretive). No isometric linear map copies the energy, $z\mapsto(z,z)$ has $E_6\oplus E_6=2E_6$, and the
genuine duplication of a quantum, $z\mapsto z\otimes z$, is quadratic; that is the no-cloning theorem \cite{WoottersZurek82} (inherited; treating the field point $z$ as a quantum state is interpretive). A
linear duplication divides the undivided energy; it does not multiply it.
\end{remark}

\begin{proposition}[An isometric complex structure exists only on the duplicate]
\label{d5:prop:cx}
\begin{enumerate}[label=(\roman*),leftmargin=2.2em]
\item No $J\in O(3,3)$ satisfies $J^2=-\id$.
\item On $\R^{6,6}$, $C_v$ satisfies $C_v^2=-\id$, $C_v^{\!\top}C_v=\id$ and $C_v^{\!\top}\eta_2C_v=\eta_2$, and commutes with $G_2$ and with the diagonal
$\lG$. It is a compact charge, invariant under the Lorentz group that acts on both frames alike.
\item $\uu6\cap\so{6,6}=\so6$ has dimension $15$; it consists of the paired rotations and is the compact part of $\Cent(G_2)$. At $n=1$ the
corresponding algebra is $\uu3\cap\so{3,3}=\so3$. Inside $\so6$ the centraliser of $C_v$ is
$\uu3'=\{\id\otimes\operatorname{diag}(A,A)+\varepsilon\otimes\operatorname{diag}(B,B):A^{\!\top}=-A,\ B^{\!\top}=B\}$, of dimension $9$, and its part
commuting with the diagonal boosts is $\R\,C_v$.
\end{enumerate}
\end{proposition}

\begin{proof}
(i) If $J^{\!\top}\eta J=\eta$ and $J^2=-\id$, then $h(x,y)=\eta(x,y)+i\,\eta(x,Jy)$ is a Hermitian form on $(\R^6,J)\cong\C^3$ with $\Re h=\eta$. A
Hermitian form of signature $(p,q)$ has real signature $(2p,2q)$, never $(3,3)$. (ii) and (iii) are computed (d5-3g, d5-3h).
\end{proof}

Part (i) is inherited, the signature parity of Hermitian forms. It bites because the arena's triples are odd, and that much is fixed. Read
physically, $w=X^{(1)}+iX^{(2)}$ combines two neutral copies into one field charged under $C_v$, as two real fields form a charged one
(inherited). The pointwise unitary symmetry $U(3)$ of $E_6$ is cut to $O(3)$ by Lorentz compatibility on one copy
(\cref{d1:sec:main}). It reappears inside $SO(6,6)$ as $U(3)'$, with the copy index as imaginary unit and not the binding root
(interpretive), and $U(3)'$ is not Lorentz invariant.

\subsection{The duplicated closure symmetry}
\label{d5:sec:closure}

\begin{proposition}[$\Cent(G_2)=\so{6,\C}$ and the copy algebra]
\label{d5:prop:so6C}
\begin{enumerate}[label=(\roman*),leftmargin=2.2em]
\item The centraliser of $G_n=G\oplus\cdots\oplus G$ in $\so{3n,3n}$ has dimension $3n(3n-1)$: $6$, $30$ and $72$ at $n=1,2,3$. It is $\so{3n,\C}$, with
compact part $\so{3n}$ of dimension $3$, $15$ and $36$. At $n=2$ every element is $\bigl(\begin{smallmatrix}R&S\\-S&R\end{smallmatrix}\bigr)$ in the
basis of the six temporal directions followed by their partners, with $R,S$ antisymmetric.
\item $\so{6,\C}=\lG^{(1)}\oplus\lG^{(2)}\oplus(\C^3\otimes\C^3)$, of dimensions $6+6+18$. Under the diagonal rotations the cross part is
$\mathbf1\oplus\mathbf3\oplus\mathbf5$ (complex), with Casimir $0$, $2$, $6$ on $2$, $6$, $10$ real dimensions.
\item The commutant of the diagonal $\lG$ in $\so{6,6}$ is the singlet of (ii), the copy algebra $\so{2,\C}=\operatorname{span}\{C_v,\varepsilon\otimes G\}$. Here
$C_v^2=-\id$ is compact, while $(\varepsilon\otimes G)^2=+\id$ is a copy boost, acting on $(z_1,z_2)$ by $\bigl(\begin{smallmatrix}0&-i\\ i&0\end{smallmatrix}\bigr)$; it
preserves $q_2$ and not $E_6\oplus E_6$.
\item $\Cent(G_2)$ meets $\lint^{(1)}$ in $0$, and the elements of $\so{6,6}$ commuting with both $G_2$ and $\lint^{(1)}$ form $\lG^{(2)}\cong\so{3,\C}$.
\end{enumerate}
\end{proposition}

Part (i) is the $n=2$ member of Theorem~17.19 of \cite{ES17}, with the same proof, and the dimension $72$ at $n=3$ re-verifies that theorem.
Part (iii) restates row three of Proposition~13.14 and Theorem~10.8 with $k_B=2$ of \cite{ES17} (fixed there). Parts (ii) and (iv) are fixed,
and (iv) is the $n=2$ analogue of Theorem~17.20 of \cite{ES17} (d5-4a to d5-4c).

\begin{proposition}[Duplicated momentum fields in the closure]
\label{d5:prop:balance}
Let one rotor act on two constituents with the same inertia $N=(\nu_1\eta^{-1}+\nu_2G\eta^{-1})/m$, whose momentum fields are a duplication of one
momentum field: $\pi^{(1)}=\alpha\pi$ and $\pi^{(2)}=\beta\pi$, with $\alpha,\beta$ acting through $\operatorname{span}\{\id,G\}$. Then
$w_1+w_2=(\alpha^2+\beta^2)\,w(\pi)$ with $w=q(z(\pi))$ (equation~(30) of \cite{ES17}). The effective strength of Proposition~17.18 of \cite{ES17}
is therefore $R_{\rm eff}=|\nu_1+i\nu_2|\,|q_2(\pi)|/m$: the closure sees the modulus of the holomorphic invariant of the duplicated momentum field. For the equal
split with a relative quarter turn of duality, $\pi^{(2)}=G\pi^{(1)}$, equivalently $p^{(2)}=-Gp^{(1)}$, one has $w_2=-w_1$ and $R_{\rm eff}=0$. The
rate is then exactly constant, and nonzero when $h(0)\neq0$, while each copy keeps $E_0(\pi^{(k)})=\tfrac12E_0(\pi)\neq0$. On one constituent with
the same inertia and $(\nu_1,\nu_2)\neq(0,0)$ a constant nonzero rate requires $E_0(\pi_0)=0$ (Corollary~17.14 of \cite{ES17}); for
$(\nu_1,\nu_2)=(0,0)$ the rate is constant for all data.
\end{proposition}

This restates Proposition~17.18 and Corollary~17.14 of \cite{ES17} (fixed there). The expression of $R_{\rm eff}$ through $q_2$ for duplicates is a
one-line consequence (fixed; d5-9f). So, for $(\nu_1,\nu_2)\neq(0,0)$, a duplicate can do what one copy cannot: run at an exactly constant rate while its copies carry nonzero
rest energy in the momentum field, in the sense that $E_0(\pi^{(k)})\neq0$. Reading $E_0(\pi^{(k)})$ as a rest energy is the interpretive dictionary stated after \cref{d1:thm:closure}. The identification $mc^2=\hbar\kappa$ concerns the common rate $\kappa=h/I_0$ of both copies, which does not depend on $E_0(\pi^{(k)})$, so it does not make $E_0(\pi^{(k)})$ a mass (interpretive).

\subsection{The five Lorentz actions at the second rung}
\label{d5:sec:actions}

What the duplication does depends on how the Lorentz algebra acts on the two copies. Theorem~10.8 of \cite{ES17} lists the choices, and at $n=2$
there are five (\cref{d5:tab:actions}).

\begin{table}[t]
\centering
\small
\begin{tabular}{@{}lccccll@{}}
\toprule
action & $(k_I,k_B,k_S)$ & algebra & dim & compact & content of $\R^{6,6}$ & spin on $S_2$\\
\midrule
spectator & $(1,0,1)$ & $\so{5,3}$ & $28$ & $13$ & $V\oplus8\cdot\mathbf1$ & half-integer\\
binding $+$ spectator & $(0,1,1)$ & $\so{3,3}$ & $15$ & $6$ & $\Lambda^2V\oplus6\cdot\mathbf1$ & integer\\
diagonal interior & $(2,0,0)$ & $\so2\oplus\so4$ & $7$ & $7$ & $2V\oplus4\cdot\mathbf1$ & integer\\
mixed & $(1,1,0)$ & $\so2$ & $1$ & $1$ & $V\oplus\Lambda^2V\oplus2\cdot\mathbf1$ & half-integer\\
binding on both & $(0,2,0)$ & $\so{2,\C}$ & $2$ & $1$ & $2\,\Lambda^2V$ & integer\\
\bottomrule
\end{tabular}
\caption{The Lorentz actions at $n=2$ and the grade-two elements commuting with them (Theorem~10.8 of \cite{ES17}), with their compact parts,
the content of the vector space (Lorentz Casimir $3$ on a four-vector $V$, $4$ on a field strength $\Lambda^2V$), and the effect of a $2\pi$
rotation on $S_2$. Every entry was computed in the realisation (d5-5a, d5-5b, d5-6b).}
\label{d5:tab:actions}
\end{table}

Read row by row: under the spectator action copy~$2$ is a passenger whose six directions join a non-compact internal algebra. With
$\lG$ on copy~$1$, copy~$2$'s whole $\so{3,3}$ becomes internal symmetry and no spin one-half survives. The diagonal interior action gives a
compact algebra containing the copy rotation $M_{12}$, but only integer spin. The mixed action keeps only $J_1=e^{(1)}_2e^{(1)}_3$. With $\lG$ on both
copies the frames are two field strengths mixed by $\so{2,\C}$. The table restates Theorem~10.8 and Propositions~10.9, 13.16 and 13.17 of
\cite{ES17} (fixed there), re-verified. For the mixed rung we also re-verified the Casimir identity
$(C_{\rm mix}-\tfrac32)(C_{\rm mix}-\tfrac{11}2)=0$ with ranks $16$ and $48$ (d5-5b). We checked the grade-three element $\Theta_F$ of Section~13.6 of
\cite{ES17}, which identifies the frame field of copy~$2$ with $dA$ of copy~$1$. It and $\Theta_F$ followed by $G$ are invariant and span the
invariants in $\Lambda^2(\text{four-plane}_1)\otimes V^{(2)}$, of dimension two as stated there. Here $\text{four-plane}_1=\langle e^{(1)}_1,e^{(1)}_4,e^{(1)}_5,e^{(1)}_6\rangle$ is the interior four-plane of copy~$1$ and $V^{(2)}$ is the span of the $e^{(2)}_a$, the space of copy~$2$ alone, which Section~13.6 of \cite{ES17} writes $V_2$. The full space of invariant grade-three elements has
dimension four: the other two lie in $\Lambda^3V^{(2)}$, among the trivectors of copy~$2$ alone, the two $\lG$-invariant trivectors of
$(\Lambda^3\R^6)_\C=2\,(0,0)\oplus2\,(1,1)$ (Section~12.3 of \cite{ES17}; d5-5c). This is a precision, not a correction. The mixed rung is the only
action at $n=2$ that houses both readings, at the price of the whole internal algebra.

\subsection{Spin and statistics of the duplicate}
\label{d5:sec:spin}

\begin{proposition}[Spin of the duplicate]
\label{d5:prop:spin}
\begin{enumerate}[label=(\roman*),leftmargin=2.2em]
\item The copy rotations $M_{12}$ (four commuting terms of square $-\tfrac14$) and $C_b$ (six) have eigenvalues $i\{-2,\dots,2\}$ and $i\{-3,\dots,3\}$, and
$\exp(2\pi M_{12})=\exp(2\pi C_b)=\id$. A surplus rotation $\tfrac12e_se_t$ has $\exp(\pi e_se_t)=-\id$.
\item A rotation through $2\pi$ is $+\id$ on $S_2$ under the diagonal interior action and $-\id$ under the spectator action. Under the mixed action it
is $-\id$, and under $\lG$ on one or both copies it is $+\id$. The spin Casimir is $0$ on $16$ and $2$ on $48$ dimensions for the diagonal action
($j=0,1$), and $\tfrac34$ on all of $S_2$ for the spectator action.
\item Under the diagonal interior action the Lorentz Casimir is $0,3,4$ on $8,32,24$ dimensions and $C_{\rm cop}=-M_{12}^2$ is $0,1,4$ on $24,32,8$, with
$C_{\rm cop}+C_\Delta=4$. The exchange $\tau_K$ commutes with the diagonal Lorentz action and anticommutes with $M_{12}$, and its $\pm1$ eigenspaces
carry identical Lorentz content, $0,3,4$ on $4,16,12$ dimensions each.
\end{enumerate}
\end{proposition}

This restates Proposition~10.9, Theorem~10.4 at $n=2$ ($\tfrac12n(n+2)=4$) and the last paragraph of Section~13.6 of \cite{ES17} (fixed there),
with the $2\pi$ eigenvalues computed (d5-6a, d5-6b). That two spin-$\tfrac12$ constituents make an integer-spin composite is inherited
\cite{EhrenfestOppenheimer31}: under the diagonal action the duplicate of a fermion is a boson, and fermions need an odd number of interior
constituents. The Lorentz isomorphism of the two exchange sectors in (iii) is fixed (d5-6c). It holds because exchange reverses the copy
charge, pairing charge $\pm m$, and splits the neutral sector evenly; there is no ``singlet antisymmetric, triplet symmetric'' rule once the
surplus directions and the Koszul signs are included.

\subsection{What the second rung adds: the vacuum and the imaginary unit}
\label{d5:sec:vacuum}

\begin{proposition}[The photon vacuum and the imaginary unit, re-verified]
\label{d5:prop:vacuum}
\begin{enumerate}[label=(\roman*),leftmargin=2.2em]
\item At $n=1$ the commutant of $\lint$ and $Q=\tfrac12e_2e_3$ in $M_8(\R)$ has dimension four, and its symmetric part is
spanned by $\id$ and $\mathcal V=e_1\cdots e_6$. The rank-four projectors in it are $\tfrac12(\id\pm\mathcal V)$, and every bivector commutes with
$\mathcal V$, so they are rigid (Proposition~26.4 of \cite{ES17}).
\item (Theorems~26.6 and~26.7 of \cite{ES17}) With $a,b=e^{(1)}_2,e^{(1)}_3$ and $c,d,e=e^{(2)}_1,e^{(2)}_2,e^{(2)}_3$: $H=abcd$ has $H^2=\id$, $H^{\!\top}=H$
and $\Tr H=0$, and $\Pi_0=\tfrac12(\id+H)$ commutes with $Q=\tfrac12ab$, with the four interior gammas and with the internal volume. Moreover
$\pi=[\tfrac12be,\Pi_0]=-\tfrac12acde$ is symmetric with $\Tr\pi^2=16$ and $[Q,\pi]=\tfrac12bcde$. The $Z_i=I_5v_i$ are symmetric of grade four with
$\{Z_i,Z_j\}=2\delta_{ij}$. $\Pi(n)$ is a rank-$32$ projector commuting with the gammas (checked at $n=(2,3,6,0,0)/7$). The tangents
$T_i=\tfrac12I_5v_i$ have $\Tr T_iT_j=16\delta_{ij}$, $[Q,T_a]=-T_b$, $[Q,T_b]=T_a$ and $[Q,T_c]=[Q,T_d]=0$.
\item The grade-four blades in the five positive internal generators that commute with $Q$ are exactly $abcd$, $abce$ and $abde$. Every
such vacuum straddles the two copies.
\item (Theorem~4.1 of \cite{ES17} at $n=2$) The product of the eight internal generators is $J=(e_2e_3)\otimes\omega$, with $J^2=-\id$. The joint commutant
of $\lint^{(1)}$ and the $28$ internal bivectors in $M_{64}(\R)$ is spanned by the blades $\id,\Omega,J,\Omega J$. The solutions of $x^2=-\id$ in it are
exactly $\pm\Omega$ and $\pm J$, and only $\pm J$ commute with the gammas.
\end{enumerate}
\end{proposition}

Parts (i), (ii) and (iv) restate results of \cite{ES17} (fixed there), re-verified in the realisation of $\Cl(6,6)$ (d5-7a to d5-7c, d5-8).
Part (iii) is fixed but elementary (d5-7d). The vacuum of Theorem~26.6 of \cite{ES17} is built on four positive internal directions, two of them the surplus pair $a,b$ of copy~$1$. A
grade-four blade in the five positive internal generators commutes with $Q$ exactly when it contains both $a$ and $b$, so it straddles the copies.
Part (iii) covers only these positive blades, not blades that also use the negative internal directions of copy~$2$. At one copy the first rung is
rigid by the commutant argument of Proposition~26.4 of \cite{ES17}, without reference to blades, so the vacuum exists because there are two copies. Part (iv) adds one
identity: the imaginary unit of the duplicate is the imaginary unit of copy~$1$ times the volume element of copy~$2$, which is that copy's
fermion parity up to sign (\cref{d5:prop:fock}). It is neither the embedded one-rung charge $J_1\otimes\id$ nor $J^{(1)}+J^{(2)}$.

\subsection{\texorpdfstring{The copy rotation as a $b$-field between two manifolds}{The copy rotation as a b-field between two manifolds}}
\label{d5:sec:bfield}

On one copy a $b$-field commuting with $G$ is removable along a worldline (\cref{d4:prop:larmor}), and a purely anticommuting one tears
the binding (\cref{d4:thm:pG}). An admissible one with both parts keeps every orbit bounded when the two triples rotate at equal rates,
outside one exceptional ratio, and it is then not removable (\cref{d4:thm:screen,d4:prop:remove}). On two copies the binding term is $K\otimes G$ with $K=\operatorname{diag}(\kappa_1,\kappa_2)$. That is a
multiple of $G_2$ only when the rates agree, and the following theorem measures what changes. The interior Lagrangian of each copy is that
of \cref{d4:sec:bfield}, with $S=K\otimes G+b$; the frame field obeys $\dot X=-SX+\pi/m$ and the momentum field $\dot\pi=S^\dagger\pi$.

\begin{theorem}[The inter-copy $b$-field]
\label{d5:thm:bfield2}
Let $\kappa_1,\kappa_2\ge0$, $\beta\in\R$, $b=\beta\,C_v$, and $M=\bigl(\begin{smallmatrix}-S&\id/m\\0&S^\dagger\end{smallmatrix}\bigr)$ on $(X,\pi)$.
\begin{enumerate}[label=(\roman*),leftmargin=2.2em]
\item Up to scale, $b$ is the only element of $\so{6,6}\cap\so{12}$ commuting with $G_2$ and the diagonal rotations. It is compact, and
$S^\dagger=K\otimes G-b$.
\item $\det(\lambda-S)=\bigl(\lambda^4+(\kappa_1^2+\kappa_2^2+2\beta^2)\lambda^2+(\kappa_1\kappa_2-\beta^2)^2\bigr)^3$, so $\spec S=\{\pm i\omega_\pm\}$ with
$\omega_\pm=\tfrac12(\kappa_1+\kappa_2)\pm\sqrt{\tfrac14(\kappa_1-\kappa_2)^2+\beta^2}$; equivalently, $\omega_\pm$ are the eigenvalues of
$\bigl(\begin{smallmatrix}\kappa_1&i\beta\\-i\beta&\kappa_2\end{smallmatrix}\bigr)$.
\item $[K\otimes G,b]=\beta(\kappa_2-\kappa_1)\,\sigma_x\otimes G$, so for $\beta\neq0$, $b$ is removable along the worldline (Larmor \cite{Larmor1897}) iff $\kappa_1=\kappa_2$. The
centraliser of $K\otimes G$ in $\so{6,6}$ has dimension $30$ if $\kappa_1=\kappa_2\neq0$, $12$ if $\kappa_1\neq\kappa_2$ are both nonzero, and $21$ if
exactly one vanishes.
\item $(M^2+\omega_+^2)(M^2+\omega_-^2)=0$ identically in $\kappa_1,\kappa_2,\beta,m$.
\item Every orbit is bounded iff either $\kappa_1=\kappa_2>0$, or $\kappa_1\neq\kappa_2$ and $\kappa_1\kappa_2\neq\beta^2$. In particular, for $\kappa_1=0<\kappa_2$
the orbits are bounded iff $\beta\neq0$: the inter-copy field confines a frame with no rate of its own. For $\kappa_1=\kappa_2=0$ they are unbounded
for every $\beta$.
\item For $\kappa_1=0<\kappa_2$, $\omega_+\omega_-=-\beta^2$ and $\omega_++\omega_-=\kappa_2$. For $\beta\ll\kappa_2$ the slow line is
$|\omega_-|=\beta^2/\kappa_2-\beta^4/\kappa_2^3+\cdots$, the seesaw. For $\beta\gg\kappa_2$ both lines approach $\pm\beta$, and the orbit size grows like
$2/(m\kappa_2)$.
\end{enumerate}
\end{theorem}

\begin{proof}
(i) is \cref{d5:prop:energysym}(i), and $C_v$ is $\eta_2$-antisymmetric while $K\otimes G$ is self-adjoint. (ii) On the sector $G_2=i$,
$S=(iK+\beta\varepsilon)\otimes\id_3$, whose characteristic polynomial is $\lambda^2-i(\kappa_1+\kappa_2)\lambda-\kappa_1\kappa_2+\beta^2$; the sector $G_2=-i$ is
its conjugate. (iii) $K\varepsilon-\varepsilon K=(\kappa_2-\kappa_1)\sigma_x$. If $\kappa_1=\kappa_2=\kappa$, then $X=e^{-b\tau}Y$ and $\pi=e^{-b\tau}\varpi$ turn
the system into $b=0$ with rate $\kappa$; if $\beta\neq0$ and $\kappa_1\neq\kappa_2$, this substitution makes the binding $e^{b\tau}(K\otimes G)e^{-b\tau}$ time-dependent. The
centraliser dimensions are computed. (iv) is a symbolic identity (d5-9c). (v) If $\kappa_1\neq\kappa_2$ and $\kappa_1\kappa_2\neq\beta^2$, then $\omega_+\neq\omega_-$,
$\omega_-\neq0$ and $\omega_++\omega_-=\kappa_1+\kappa_2>0$, so the polynomial of (iv) has four simple imaginary roots and $M$ is semisimple with imaginary
spectrum. If $\kappa_1=\kappa_2=\kappa>0$, use (iii) and the boundedness of the $b=0$ system (Theorem~16.2 of \cite{ES17}); if $\kappa=0$ the substituted
system is free drift. If $\kappa_1\neq\kappa_2$ and $\kappa_1\kappa_2=\beta^2$, then $\omega_-=0$, and $M$ is semisimple at $0$ iff $\ker S^\dagger\subset\operatorname{im}S$.
Because $S^{\!\top}=-S$ and $S^\dagger=-\eta_2S\eta_2$, this says that $\ker S$ is $\eta_2$-totally isotropic. Complexifying, $\ker S$ is spanned by
the real parts of $v\otimes c$ in the $G_2=i$ sector, with $c\in\C^3$ in the basis $e_j-ie_{j+3}$ and $v=(\beta,i\kappa_1)$ (or $v=(1,0)$ when $\beta=\kappa_1=0$), and
$\eta_2(\Re(v\otimes c),\Re(v\otimes c'))=\Re\bigl[(v\cdot v)(c\cdot c')\bigr]$. This vanishes for all $c,c'$ iff $v\cdot v=\beta^2-\kappa_1^2=0$, which together
with $\kappa_1\kappa_2=\beta^2$ means $\kappa_1=\kappa_2$ or $\beta=\kappa_1=0$, and in the second case $v=(1,0)$ is not isotropic; so there is a Jordan block.
(vi) Vieta and expansion. The orbit sizes are numerical (d5-9e, d5-9g).
\end{proof}

Parts (ii) and the seesaw in (vi) are inherited: two-level mixing with detuning $\tfrac12(\kappa_1-\kappa_2)$ and coupling $\beta$, the generalised Rabi
splitting, and the seesaw formula \cite{Minkowski77}. Part (i) restates Proposition~13.14 of \cite{ES17} (fixed there). Parts (iii)--(v) are fixed
(d5-9a to d5-9d). The exact checks are: bounded at $(\kappa_1,\kappa_2,\beta)=(1,3,\tfrac25),(0,5,\tfrac12),(2,2,\tfrac13),(1,1,1),(2,2,0)$ and
$(0,\tfrac1{100},1)$; unbounded at $(0,0,\tfrac12),(1,4,2),(1,9,3),(0,3,0)$. At $(0,5,\tfrac12)$, $\omega_-=-0.04951$ against $\beta^2/\kappa_2=0.05$, and
$\sup_{t\le600}\|e^{Mt}\|=19.86$ against $1/|\omega_-|=20.2$. At $(0,\tfrac1{100},1)$ the lines are $1.00501$ and $-0.99501$, and the supremum over a beat
period is $200.0=2/(m\kappa_2)$. So the seesaw orbit is bounded, with amplitude of order $1/(m|\omega_-|)$ for $\beta\ll\kappa_2$ and $2/(m\kappa_2)$ for
$\beta\gg\kappa_2$.

\begin{remark}[The author's first remark at two manifolds]
\label{d5:rem:dealbreaker}
On one manifold a $b$ commuting with $G$ is removable along a worldline, and a constant connection alone gives free drift
(\cref{d4:prop:remove,d4:prop:alone}). An admissible $b$ that does not commute with $G$ can already there be non-removable, change the spectrum
and keep every orbit bounded (\cref{d4:thm:screen}), so these two features are not new at $n=2$. What two manifolds with different rates add is
this. The compact inter-copy field commutes with $G_2$, so on one manifold a field of that kind would be removable, but it does not commute with
the binding term $K\otimes G$. It is therefore not removable, it shifts the lines, and it confines a frame with $\kappa_1=0$ by coupling it to a
frame with $\kappa_2>0$. The no-go of \cref{d4:prop:alone} still holds at $n=2$: for $\kappa_1=\kappa_2=0$ the motion is free drift for every
$\beta$. In that sense, and only in it, the $b$-field is a ``deal breaker'' at $n\ge2$ (fixed). Two qualifications remain. The field can set the lines, which approach $\pm\beta$ when $\beta\gg\kappa_2$, but not the rates, and the size of
the orbit is set by $\kappa_2$ and diverges as $\kappa_2\to0$. And reading a line or a rate as a rest energy still needs $mc^2=\hbar\kappa$ (an added
identification).
\end{remark}

\subsection{Assessment}
\label{d5:sec:assess}

\paragraph{What is forced, what is chosen, what is added.} \emph{Forced} (inherited): the algebra step, $\Cl(V\oplus V)=\Cl(V)\hat\otimes\Cl(V)$, with the
Klein factor determined up to conjugation; $S_2=S\otimes S$ as the Fock space of six modes; and the exchange as the Koszul swap.
\emph{Chosen} (constructed): the Lorentz action, one of five; the energy duplication, a point of the three-sphere, or of the circle if $\eta$ is
to be respected, with the angle $\phi$ free; and the vacuum point on $S^4$. \emph{Added} at $n=2$ (fixed; the copy algebra, the copy charge, the constant-rate closure, the vacuum and the mixed rung restate \cite{ES17}, while the complex-structure statement of \cref{d5:prop:cx}, the comparison $\so6\supset\uu3'$, the form of $R_{\rm eff}$ through $q_2$ and the inter-copy $b$-field are fixed here): the copy algebra $\so{2,\C}$ on two
binding frames (dimension $0$ at $n=1$); an $\eta$-isometric complex structure and a compact, Lorentz-invariant copy charge (none at $n=1$); a
Lorentz-compatible energy symmetry $\so6\supset\uu3'$, against $\so3$; a closure with constant rate and non-null copies, $E_0(\pi^{(k)})\neq0$; a non-rigid,
Lorentz-invariant vacuum; the mixed rung, which houses both readings; and an inter-copy $b$-field that confines (d5-10 collects the
dimensions). \emph{Lost}: spin one-half under the diagonal interior action, and the uniqueness of the photon direction (Corollary~6.6 of
\cite{ES17}).

\paragraph{Processes.} \Cref{d5:tab:processes} lists the physical processes that take one copy to two, and what each preserves.

\begin{table}[t]
\centering
\footnotesize
\begin{tabular}{@{}>{\raggedright\arraybackslash}p{3.3cm}>{\raggedright\arraybackslash}p{4.7cm}>{\raggedright\arraybackslash}p{4.3cm}>{\raggedright\arraybackslash}p{2.4cm}@{}}
\toprule
process & mathematics & preserves & class\\
\midrule
splitting one field into two (beam splitter) & $\iota_{\alpha,\beta}$, $S^3=U(2)/U(1)$; circle if also $\eta$-isometric & $E_6$, $T$, $\sum E_0^{(k)}$; $q_2$ only on the circle & inherited; reading fixed\\
second quantisation & $\Cl(V\oplus V)=\Cl(V)\hat\otimes\Cl(V)$; Fock space; Jordan--Wigner strings & Clifford relations & inherited (forced)\\
two neutral copies into one charged field & $C_v$ as $i$, $w=X^{(1)}+iX^{(2)}$ & $\eta_2$, $E_6\oplus E_6$ & inherited; exists only at $n\ge2$: fixed\\
two-body composite & diagonal interior action; exchange $\tau_K$ & integer spin & inherited; restated\\
pair creation from null energy & linear covariant duplication keeps null copies; two non-collinear null fields have $E_{\rm pair}>0$ & $T$ & inherited; pair-creation reading interpretive\\
cloning, $z\mapsto z\otimes z$ & quadratic, not isometric and linear & --- & obstructed (inherited)\\
closure balance & $\pi^{(2)}=G\pi^{(1)}$, $R_{\rm eff}\propto|q_2(\pi)|=0$ & constant rate, $E_0^{(k)}\neq0$ & restated (Prop.~17.18)\\
inter-copy $b$-field & $S=K\otimes G+\beta C_v$ & bounded iff \cref{d5:thm:bfield2}(v) & fixed; mixing inherited\\
\bottomrule
\end{tabular}
\caption{Processes that take one copy of the arena to two.}
\label{d5:tab:processes}
\end{table}

\paragraph{The precise sense of ``a unitary energy results in a duplication''.} Three statements, in increasing order of what the arena adds.
First, the undivided energy $(\C^3,E_6,U(3))$ embeds isometrically in $(\C^3\otimes\C^2,E_6\oplus E_6,U(6))$. The Lorentz-covariant embeddings
form the three-sphere of a beam splitter, and the diagonal $U(3)$ and the copy $U(2)$ are a dual pair (inherited). Second, compatibility with
$\eta$ cuts the copy $U(2)$ to $O(2)$ and $U(6)$ to $O(6)$, since $U(n)\cap O(n,\C)=O(n)$. The copy rotation $SO(2)$, the identity component of
$O(2)$, is the one connected family that preserves the energy, the metric, the binding root and the diagonal Lorentz group at once; the other
component contains the swap of the copies and $\operatorname{diag}(1,-1)$. Neither exists on one copy (fixed). Third, the duplicate's rest
energy splits into a stress-energy part, preserved by every duplication, and a holomorphic part. The complex invariant $q_2$ is preserved only on the circle of the copy rotation. Its modulus $|q_2|$, which is what
the closure sees, is preserved on that circle composed with the overall duality rotation $G_2$, which multiplies $q_2$ by a phase (fixed). Read literally as a multiplication of the energy, the phrase fails: a linear duplication divides the
undivided energy, and copying it is impossible (no-cloning, inherited).

\begin{table}[t]
\centering
\footnotesize
\begin{tabular}{@{}>{\raggedright\arraybackslash}p{8.1cm}>{\raggedright\arraybackslash}p{5.0cm}>{\raggedright\arraybackslash}p{2.0cm}@{}}
\toprule
statement & class & checks\\
\midrule
\cref{d5:prop:graded}: graded tensor product, $\Cl(6,6)=M_{64}(\R)$, Klein factor forced & inherited (restated: eq.~(2)) & d5-1a to d5-1d\\
\cref{d5:prop:fock}: Fock space, Jordan--Wigner strings, pole operators $=\id-2N_i$ & inherited (reading of Prop.~2.5) & d5-1e\\
\cref{d5:prop:cmrel}: halves, $\pi/4$ copy rotation, $\tau_K=g_1\cdots g_6$, what each half carries & inherited; identities fixed & d5-1f, d5-2a to d5-2c\\
\cref{d5:prop:energysym}: $\uu2$, $\so{2,\C}$, $\so2$; dual pair & fixed (restated: Prop.~13.14); inherited & d5-3a, d5-3b\\
\cref{d5:thm:dup}: duplication maps, $S^3$ and circle, $q_2$ versus $T$ & inherited; reading fixed & d5-3c to d5-3e\\
\cref{d5:rem:twoE0}: two rest energies; Gram invariants; pair invariant; no cloning & fixed; inherited; analogies interpretive & d5-3f, d5-3i\\
\cref{d5:prop:cx}: isometric complex structure only at $n\ge2$; $\so6\supset\uu3'$ & inherited (parity); fixed & d5-3g, d5-3h\\
\cref{d5:prop:so6C}: $\Cent(G_n)=\so{3n,\C}$; $\mathbf1\oplus\mathbf3\oplus\mathbf5$; copy algebra; analogue of Thm~17.20 & fixed (restated: Thm~17.19, Prop.~13.14) & d5-4a to d5-4c\\
\cref{d5:prop:balance}: closure sees $|q_2(\pi)|$; balance with non-null copies & fixed (restated: Prop.~17.18, Cor.~17.14) & d5-9f\\
\cref{d5:tab:actions}: five Lorentz actions; mixed rung; $\Theta_F$ (dimension $2$ within $\Lambda^2(\text{four-plane}_1)\otimes V^{(2)}$, $4$ in all) & fixed (restated: Thm~10.8, Props.~13.16--13.17) & d5-5a to d5-5c\\
\cref{d5:prop:spin}: spin and statistics, Casimirs, exchange sectors & fixed (restated: Prop.~10.9, Thm~10.4); inherited & d5-6a to d5-6c\\
\cref{d5:prop:vacuum}: rigid first rung, $S^4$ vacuum, straddling, $J=J_1\otimes\omega$ & fixed (restated: Prop.~26.4, Thms~26.6, 26.7, 4.1) & d5-7a to d5-7d, d5-8\\
\cref{d5:thm:bfield2}: inter-copy $b$-field: uniqueness, spectrum, removability, boundedness criterion, seesaw & fixed; mixing and seesaw inherited & d5-9a to d5-9e, d5-9g\\
\cref{d5:rem:dealbreaker}: the deal-breaker reading & fixed; obstructed for $\kappa_1=\kappa_2=0$; $mc^2=\hbar\kappa$ an identification & ---\\
\cref{d5:sec:assess}: comparison $n=1$ against $n=2$ & fixed; mostly restated (Prop.~13.14) & d5-10\\
\bottomrule
\end{tabular}
\caption{Claims of \cref{d5:sec:main}, their classes, and the checks printed by \texttt{d5\_duplication\_math.py}. ``Restated'': fixed in
\cite{ES17}, restated and re-verified here.}
\label{d5:tab:claims}
\end{table}
